\section{Introduction}
The rapid advancement of Large Language Models (LLMs) \citep{qwen2,llama3,gpt4o,deepseekv3} has marked a pivotal stage in artificial intelligence, achieving state-of-the-art performance across a wide range of tasks. Recently, a new class of \textbf{reasoning models} has emerged to explicitly tackle complex reasoning problems \citep{deepseekr1,openaio1,kimik1.5,qwen3,llama4}. These models rely on long Chain-of-Thought (CoT) reasoning and generate extensive intermediate steps during inference. While highly effective in domains such as mathematical reasoning and code generation, they introduce new challenges for efficient training and deployment.

A key drawback of reasoning models is their tendency to generate a large number of output tokens \citep{overthinking,overthinking1,speculativecot}, which substantially inflates the KV cache and leads to severe memory and bandwidth bottlenecks during attention computation. Although prior work has explored KV cache compression, most methods are designed for long-input, short-output scenarios and are ill-suited for long-output generation. Some approaches only compress the KV cache during the prefill stage \citep{snapkv,pyramidkv,calidrop}, while others incur significant computational overhead and require additional auxiliary storage \citep{h2o,vatp}. These limitations are particularly pronounced in long-output settings, highlighting the need for more efficient KV cache compression methods for reasoning models.

To address this gap, we propose \textbf{LongFlow}, an efficient KV cache compression method tailored for long-output generation. The core idea of LongFlow is a fast yet accurate importance estimation metric for historical tokens, derived directly from an intermediate value of the standard attention computation using only the current query. This design requires no auxiliary storage and introduces negligible computational overhead. We provide both theoretical justification, showing that the metric approximates attention output loss, and empirical validation through extensive experiments.

Beyond the algorithmic design, we introduce several system-level optimizations to maximize practical efficiency. We adopt a static KV cache that pre-allocates memory to avoid fragmentation and dynamic allocation overhead. To further improve hardware utilization, we implement a custom Triton kernel that fuses FlashAttention, importance estimation, and token eviction into a single optimized operator, reducing attention latency from 47\,ms to 8\,ms (Figure~\ref{fig:speed_kernel}). Overall, LongFlow achieves an 11.8$\times$ throughput improvement with an 80\% KV cache compression ratio, while preserving most of the model accuracy.

Our contributions are summarized as follows:
\begin{itemize}[itemsep=3pt,topsep=0pt,parsep=0pt]
\item \textbf{A lightweight KV cache compression algorithm for long-output generation.} We propose LongFlow, which introduces an accurate importance metric computed efficiently from the current query and intermediate attention results, with negligible overhead.
\item \textbf{A fused, high-performance attention kernel.} We design a custom Triton kernel that integrates attention computation, importance estimation, and token eviction into a single optimized operator.
\item \textbf{State-of-the-art efficiency for reasoning models.} Experiments demonstrate up to an 11.8$\times$ throughput improvement and an 80\% reduction in KV cache size, with minimal impact on model accuracy.
\end{itemize}
